\section{安全性与开销分析}

\subsection{不可篡改性与可追溯性}

TrustFlowAudit 通过两类哈希结构提供证据可追溯性。第一，分片哈希构成 Merkle 根，链上保存根哈希后，任一分片修改都会导致叶子哈希和路径验证失败。第二，审计事件在链下按哈希链组织，链上周期性锚定哈希链状态或批量 Merkle 根，能够发现链下日志被插入、删除或修改的情况。由于联盟链记录具有共识确认，已经提交的审计摘要、恢复状态和信誉变化难以被单方篡改。

\subsection{审计委员会的安全边界}

委员会机制降低了中低风险事件的确认开销，但也引入委员会选择风险。如果攻击者能够长期控制委员会多数节点，则可能伪造审计摘要或拒绝确认异常事件。因此，委员会选择必须满足三个条件：其一，候选节点需要经过可信评分过滤；其二，委员会需要周期性重选，避免固定节点长期垄断；其三，高风险数据对象必须提高委员会规模或回退 BLoP 式全可信节点复核。本文方法适合联盟链准入明确、节点身份可追溯的场景，不适合开放匿名网络。

\subsection{链上开销}

链上开销主要来自审计 payload 数量、payload 大小和共识确认延迟。分级批量存证将低风险事件从逐条上链转化为批量摘要上链，可以减少链上记录数；可信委员会减少每次确认参与节点数量，可以降低 PBFT 风格投票的通信压力；纠删码证据存储降低链下证据冗余，相比多副本策略更适合大规模审计日志保存。

需要注意的是，本文当前分析是方法层面的定性分析。不同联盟链平台、节点规模、网络延迟和智能合约实现会显著影响实际性能。因此，最终论文仍需要在统一实验环境下比较 BLoP 式全可信节点 PBFT、Merkle 批量存证、可信委员会和 TrustFlowAudit 全量优化四类策略。其中 BLoP 式全可信节点 PBFT 应作为强安全基线。

\subsection{与 PDP/PoR 的关系}

TrustFlowAudit 不替代 PDP/PoR，而是为其提供链上治理和闭环执行框架。PDP/PoR 负责证明远程数据持有性或可恢复性，TrustFlowAudit 负责把挑战策略、证明摘要、异常处理、恢复状态和责任记录组织为可追溯流程。当前原型用哈希抽样代替 PDP/PoR，是为了先验证系统流程；后续可以将哈希校验模块替换为公开可验证的 PDP/PoR 证明，从而提高密码学强度。
